\documentclass[11pt]{article}
\usepackage{../homework}

\profname{Dr. David Singer}
\assignnum{Assignment 04}
\duedate{28 September 2026}
\classcode{MATH 465 Differential Geometry}

\begin{document}
All problems from Do Carmo, \textit{Differential Geometry of Curves and Surfaces} (2e).

\begin{problem}{2.4.7}
	Let \( f: S \rightarrow \mathbb{R} \) be given by \( f(p)= |p-p_0|^2 \), where \( p \in S \) and \( p_0 \) is a fixed point of \( \mathbb{R}^3 \) (see example 1 of 2-3).
	Show that \( d f_p(w) = 2w \cdot (p-p_0) \), \( w \in T_pS \).
\end{problem}
\begin{proof}
	Let \(\alpha: I \rightarrow S\) be a curve on the surface such that \(\alpha(0)=p\) and \(\alpha^{\prime}(0)=w\).
	Then by definition of the differential,
	\begin{equation}
		\label{eq:5}
		\begin{aligned}
			df_p(w) & = (f \circ \alpha)^{\prime}(0) = \frac{\mathrm{d} }{\mathrm{d} t}\left( \left| \alpha(t) - p_0 \right|^2 \right) \eval_{t=0} = 2 \left| \alpha(t) - p_0 \right| \frac{\mathrm{d} }{\mathrm{d} t}|\alpha(t)-p_0| \eval_{t=0} \\
			        & = 2 \left| p-p_0 \right| \frac{2 \langle \alpha^{\prime}(t), \alpha(t) - p_{0} \rangle}{2 |\alpha(t) - p_{0}|} \eval_{t=0} = 2 \langle w, p - p_{0} \rangle.
		\end{aligned}
	\end{equation}
\end{proof}

\begin{problem}{2.4.8}
	Prove that if \( L: \mathbb{R}^3 \rightarrow \mathbb{R}^3 \) is a linear map and \( S \subset \mathbb{R}^3 \) is a regular surface invariant under \( L \), i.e., \( L(S) \subset S \), then the restriction \( L|_S \) is a differentiable map and
	\begin{equation}
		\label{eq:1}
		dL_p(w) = L(w), \quad p \in S, w \in T_pS.
	\end{equation}
\end{problem}
\begin{proof}
  For differentiability, we know that all linear maps are differentiable in \( \mathbb{R}^3 \).
  Moreover, since \( L(S) \subset S \), the restriction \( L|_S : S \rightarrow S \) is a differentiable map restricted to a regular surface and is hence differentiable.
  Let \( \alpha(s): I \rightarrow S \) be some curve on the surface for which \( \alpha(0) = p \), \(\alpha^{\prime}(0)= w\).
  Then, by definition,
  \begin{equation}
  \label{eq:19}
  \begin{aligned}
    dL_p(w) &= \frac{\mathrm{d} }{\mathrm{d} t}(L \circ \alpha)(0) = \lim_{h \rightarrow 0}\frac{L(\alpha(t + h)) - L(\alpha(t))}{h} = \lim_{h\rightarrow0} \frac{L(\alpha(t+h) - \alpha(t))}{h} \\
    &= L(\lim_{h\rightarrow0} \frac{\alpha(t+h) - \alpha(t)}{h}) = L(\alpha^{\prime}(t))\eval_{t=0} = L(w),
  \end{aligned}
\end{equation}
where we have used linearity in the first line, and the continuity of \( L \) to bring the limit inside in the second line.
\end{proof}

\begin{problem}{2.4.10 (Tubular Surfaces)}
	Let \(\alpha: I \rightarrow \mathbb{R}^3\) be a regular parametrized curve with nonzero curvature everywhere and arc length as parameter.
	Let
	\begin{equation}
		\label{eq:2}
		\mathbf{x}(s,v) = \alpha(s) + r(N(s)\cos v + B(s)\sin v), \quad r = \text{const.} \neq 0, s \in I,
	\end{equation}
	be a parametrized surface (the \textit{tube} of radius \( r \) around \(\alpha\)), where \( N \) is the normal vector and \( B \) is the binormal vector of \(\alpha\).
	Show that, when \( \mathbf{x} \) is regular, its unit normal vector is
	\begin{equation}
		\label{eq:3}
		\nu(s,v) = -(N(s)\cos v + B(s) \sin v).
	\end{equation}
\end{problem}
\begin{proof}
  By definition, \(\nu = \frac{\mathbf{x}_{s} \wedge \mathbf{x}_v}{|\mathbf{x}_s \wedge \mathbf{x}_v |}\).
  We first calculuate \( \mathbf{x}_{s} \) and \( \mathbf{x}_{v} \).
  \begin{equation}
  \label{eq:20}
  \begin{aligned}
    &\mathbf{x}_s = \alpha^{\prime}(s) + r \left( N^{\prime}(s)\cos v + B^{\prime}(s) \sin v \right), \\
    & \mathbf{x}_v = r \left( -N(s)\sin v + B(s)\cos v \right).
  \end{aligned}
  \end{equation}

  Now, the outer product is linear, so that
  \begin{equation}
  \label{eq:21}
  \begin{aligned}
    \mathbf{x}_s \wedge \mathbf{x}_v &= -r (\alpha^{\prime}(s) \wedge N(s)\sin v) + r( \alpha^{\prime}(s) \wedge B(s)\cos v ) \\
                                &\quad+ r^2 (-N^{\prime}(s)\cos v \wedge N(s) \sin v + N^{\prime}(s)\cos v \wedge B(s) \cos v \\
                                &\quad- B^{\prime}(s) \sin v \wedge N(s)\sin v + B^{\prime}(s) \sin v \wedge B(s) \cos v).
  \end{aligned}
  \end{equation}

  Recalling the Frenet equations,
  \begin{equation}
  \label{eq:22}
  \begin{aligned}
    T^{\prime} &= \kappa N \\
    N^{\prime} &= -\kappa T - \tau B \\
    B^{\prime} &= \tau N,
  \end{aligned}
  \end{equation}

  we see that \( B^{\prime} \) and \( N \) are colinear, so any terms involving the outer product of both vanish.
  Moreover, since \( \alpha \) is arc-length parametrized, \(\alpha^{\prime} = T\).
  \cref{eq:21} becomes
  \begin{equation}
  \label{eq:23}
  \begin{aligned}
    \mathbf{x}_s \wedge \mathbf{x}_v &= r(T \wedge B \cos v) -r (T \wedge N \sin v) + r^2(N^{\prime} \cos v \wedge N \sin v \\
                                &\quad+  N^{\prime}\cos v \wedge B(s) \cos v + B^{\prime}(s) \sin v \wedge B(s) \cos v) \\
    &= r(-N\cosv) - rB\sin v + r^2\cos v \sin v(-B-T) + r^2
  \end{aligned}
  \end{equation}
\end{proof}

\begin{problem}{2.5.1}
	Compute the first fundamental forms of the following parametrized surfaces where they are regular:
	\begin{enumerate}[label=(\alph*)]
		\item \( \mathbf{x}(u,v) = (a \sin u \cos v, b \sin u \sin v, c \cos u) \); ellipsoid.
		      \begin{answer}
			      \( I_{ij} = \langle \mathbf{x}_i, \mathbf{x}_j  \rangle \).
			      Hence,
			      \begin{equation}
				      \label{eq:6}
				      \begin{aligned}
					      I_{11} & = \langle (a\cos u \cos v, b \cos u \sin v, -c\sin u), (a\cos u \cos v, b \cos u \sin v, -c\sin u) \rangle \\
					             & = a^2 \cos^2 u \cos^2 v + b^2 \cos^2 u \sin^2 v + c^2 \sin^2 u.
				      \end{aligned}
			      \end{equation}

			      \begin{equation}
				      \label{eq:7}
				      \begin{aligned}
					      I_{22} & = a^2 \sin^2 u \sin^2 v + b^2 \sin^2 u \cos^2 v.
				      \end{aligned}
			      \end{equation}

			      \begin{equation}
				      \label{eq:8}
				      \begin{aligned}
					      I_{12} & = I_{21} = a^2(-\cos u \sin u \cos v \sin v) + b^2(\cos u \sin v \sin u \cos v) \\
					             & = (b^2 - a^2)(\cos u \cos v \sin u \sin v).
				      \end{aligned}
			      \end{equation}
		      \end{answer}
		\item \( \mathbf{x}(u,v) = (au \cos v, bu\sin v, u^2) \); elliptic paraboloid.
		      \begin{answer}
			      \begin{equation}
				      \label{eq:9}
				      \begin{aligned}
					      I_{11} = a^2 \cos^2 v + b^2 \sin^2 v + 4 u^2.
				      \end{aligned}
			      \end{equation}
			      \begin{equation}
				      \label{eq:10}
				      \begin{aligned}
					      I_{22} = a^2 u^2 \sin^2 v + b^2 u^2 \cos^2 v.
				      \end{aligned}
			      \end{equation}

			      \begin{equation}
				      \label{eq:11}
				      \begin{aligned}
					      I_{12} & = I_{21} = a^2(-u\cos v\sin v) + b^2(u\sin v \cos v) \\
					             & = (b^2 - a^2)(u \cos v \sin v).
				      \end{aligned}
			      \end{equation}
		      \end{answer}
		\item \( \mathbf{x}(u,v) = (au \cosh v, bu \sinh v, u^2) \); hyperbolic paraboloid.
		      \begin{answer}
			      \begin{equation}
				      \label{eq:12}
				      \begin{aligned}
					      I_{11} & = a^2 \cosh^2 v + b^2 \sinh^2 v + 4 u^2. \\
					      I_{22} & = a^2 u^2\sinh^2 v + b^2 u^2 \cosh^2 v.
				      \end{aligned}
			      \end{equation}

            \begin{equation}
            \label{eq:13}
            \begin{aligned}
              I_{12} &= I_{21} = a^2(u \cosh v \sinh v) + b^2(u \sinh v \cosh v) \\
              &= (a^2 + b^2)(u \cosh v \sinh v).
            \end{aligned}
            \end{equation}
		      \end{answer}
		\item \( \mathbf{x}(u,v) = (a \sinh u \cos v, b \sinh u \sin v, c \cosh u) \); hyperboloid of two sheets.
          \begin{answer}
            \begin{equation}
            \label{eq:15}
            \begin{aligned}
              I_{11} &= a^2 \cosh^2u \cos^2 v + b^2 \cosh^2 u \sin^2 v + c^2 \sinh^2 u. \\
              I_{22} &= a^2 \sinh^2 u \sin^2 v + b^2 \sinh^2 u \cos^2 v.
            \end{aligned}
            \end{equation}

            \begin{equation}
            \label{eq:16}
            \begin{aligned}
              I_{12} &= I_{21} = a^2(-\cosh u \sinh u \cos v \sin v) + b^2(\cosh u \sinh u \cos v \sin v) \\
              &= (b^2 - a^2)(\cosh u \sinh u \cos v \sin v).
            \end{aligned}
            \end{equation}
          \end{answer}
	\end{enumerate}
\end{problem}

\begin{problem}{2.5.9}
	Show that a surface of revolution can always be parametrized so that
	\begin{equation}
		\label{eq:4}
		E = E(v), \quad F = 0, \quad G = 1.
	\end{equation}
\end{problem}
\end{document}
